% results.tex — Results section content
% Included from main.tex via % results.tex — Results section content
% Included from main.tex via % results.tex — Results section content
% Included from main.tex via % results.tex — Results section content
% Included from main.tex via \input{sections/results}

\section{Results}
\label{sec:results}

\input{figures/tab_results}

% ============================================================
% 3.1 THE TRANSITION IS UNIVERSAL BUT MODEL-DEPENDENT
% ============================================================
\subsection{The transition is universal but model-dependent}
\label{sec:transition}

Every model we tested eventually shows greatly reduced instruction-following as the number of hardcoded turns~$N$ increases, in at least some of the experimental conditions, confirming that induction pressure is a universal phenomenon. However, the transition point varies considerably across models (Table~\ref{tab:results}; Figure ~\ref{fig:heatmaps}; individual model results in Appendix~\ref{app:permodel}). In \emph{fixed-output} conditions, average instruction-following (IF) rates across all $N$ values range from 99\% (Llama~3.3~70B) to 1\% (Qwen3~235B). At $N{=}1$, nearly all models follow instructions at high rates, but by $N{=}3$ only six of thirteen models remain above 50\%. The transition boundary---the first $N$ at which IF rate drops to or below 50\%---spans from $N{=}1$ (GPT-5.2, Hermes-4~70B, Gemini~2.5~Flash, Qwen3~235B) to never crossing the threshold within our range for Llama~3.3~70B, which maintains near-perfect instruction-following at every $N$ tested, including $N{=}50$.

\paragraph{Ranking is consistent across condition families.}
The model ranking is broadly preserved between \emph{fixed-output} and \emph{task-based} conditions, though \emph{task-based} conditions, which require the model to generate full responses rather than single tokens, show higher average IF rates overall. In these conditions, GPT-5.2 maintains the latest transition point ($N{=}43$), followed by Claude~Opus~4.6 ($N{=}31$), then Claude~Sonnet~4.6 and Llama~3.3~70B ($N{=}22$).

\subsection{Robustness is easier for task-based instructions.}
\label{sec:diversity}

\emph{Task-based} conditions consistently produce higher instruction-following rates than \emph{fixed-output} conditions (Figure \ref{fig:heatmap-dynamic}; Table~\ref{tab:results}). We show in \S\ref{sec:diversity-followup} that this gap is driven by \emph{output diversity} rather than engagement with the question content.

\paragraph{Capability predicts task-based but not fixed-output robustness.}
We correlate each model's average IF rate with two external benchmarks: GPQA \citep{rein_gpqa_2023} and IFBench \citep{pyatkin_generalizing_2025} (Figure~\ref{fig:capability}; additional analysis in Appendix~\ref{app:capability}). Once again, the results differ between \emph{fixed-output} and \emph{task-based} conditions. In \emph{fixed-output} conditions, neither benchmark achieves a significant correlation (all $p > 0.28$). In \emph{task-based} conditions, both correlations are positive, with IFBench approaching significance ($r{=}0.55$, $p{=}0.052$) and GPQA showing a moderate trend ($r{=}0.35$, $p{=}0.24$). This divergence suggests that \emph{task-based} robustness draws on some of the same capabilities measured by general benchmarks, whereas \emph{fixed-output} robustness is a distinct property unrelated to general intelligence.

\input{figures/fig_heatmaps}
% \input{figures/fig_capability}

\paragraph{The Llama anomaly.}
Both Llama~3.1 and 3.3~70B perform substantially above the next best model in \emph{fixed-output} conditions. Even at $N{=}50$, Llama~3.3 maintains an IF rate of 95\%. This outlier behavior persists across all condition types and is not explained by model scale or intelligence (both are 70B parameters, smaller than the largest models tested). Possible explanations include Meta's emphasis on instruction adherence during post-training, chat template effects, or training data composition. We consider this an open question warranting further investigation. Appendix \ref{app:training} shows further comparison between the two Llama versions.


% ============================================================
% 3.2 OUTPUT DIVERSITY EXPLAINS TASK-BASED ROBUSTNESS
% ============================================================

% ============================================================
% 3.3 CONTENT ALIGNMENT MODULATES THE TRANSITION
% ============================================================
\subsection{Content alignment modulates the transition}
\label{sec:alignment}

\input{figures/tab_alignment_gap}

We find that the content of the instruction modulates how strongly models adhere to it under induction pressure (Table~\ref{tab:alignment-gap}, Figure~\ref{fig:alignment-curves}; per-model plot in Figure~\ref{fig:alignment-gap}, Appendix~\ref{app:percondition}). In \emph{fixed-output} conditions, \emph{value-aligned} instructions (such as ``always respond with `I love being helpful'\,'') yield a mean IF rate of 32\%, compared to 18\% for \emph{value-misaligned} instructions (``always respond with `I hate being helpful'\,''), a gap of 14 percentage points. We also test a factual alignment condition which shows a similar pattern: factually correct instructions lead to higher IF rates than factually false ones (32\% vs 22\%).

The alignment effect is highly uneven across models. Claude~Opus~4.6 shows the largest gap: 87\% aligned vs.\ 6\% misaligned (an 81-point difference). Claude~Sonnet~4.6 (28-point gap) and GPT-5.2 (22-point gap) also show large effects. In contrast, Llama~3.3~70B (1-point gap) and Gemma-3~27B (0-point gap) are effectively alignment-neutral, either because they follow instructions regardless (Llama) or because they do not (Gemma).


\subsection{Output diversity, not task engagement, explains the gap}
\label{sec:diversity-followup}

Two properties distinguish \emph{task-based} from \emph{fixed-output} conditions in our experiment: \emph{task-based} conditions require the model to engage with the question content to produce a well-formed response, and their outputs are varied and natural-sounding across turns, whereas \emph{fixed-output} conditions require only a single repeated token. Either property, or both, could explain the robustness gap, so we design two follow-up conditions to disentangle them. In the \emph{classify-question} condition the model reads each question and classifies it (single-token output, but question-engaged); the \emph{random-facts} condition produces 1--3 sentence free-text responses about a fixed topic, unrelated to the input question (longer, diverse output, question-ignored). Evaluated on all 13 core models (Table~\ref{tab:followup-main}, Figure~\ref{fig:followup-curves}; per-model breakdown in Appendix~\ref{app:followup}), content engagement is not the primary driver: \emph{classify} gives only a modest improvement over the \emph{fixed-output} baseline ($+7$ percentage points, $p{=}0.034$) and stays below \emph{task-based}, while \emph{random-facts} matches or slightly exceeds \emph{task-based} robustness ($+7$ points vs.\ \emph{task-based}, $p{=}0.10$). Output diversity, not semantic engagement with the input, is therefore the primary factor distinguishing robust from fragile conditions; see \S\ref{sec:discussion} for interpretation.

\input{figures/tab_followup_main}



% ============================================================
% 3.4 TRAINING METHOD MATTERS MORE THAN SCALE
% ============================================================
% \subsection{Training method matters more than scale}
% \label{sec:training}

% We exploit two natural comparisons to isolate the contribution of training method: OLMo~3.1~32B at three post-training stages, and Llama~3.1 vs.\ 3.3~70B.

% \paragraph{OLMo training stages.}
% In fixed-output conditions, OLMo~3.1~32B-SFT (supervised finetuning only) achieves an average IF rate of 0.04, collapsing to near-zero by $N{=}3$. Adding DPO triples the average to 0.11, with IF rates of 0.85 at $N{=}1$ (vs.\ 0.43 for SFT-only). However, adding RLVR on top of DPO yields no further improvement: SFT+DPO+RLVR achieves an average IF of 0.11, statistically indistinguishable from SFT+DPO. All three stages collapse to zero by $N{=}16$. This result 

% The pattern is amplified in task-based conditions (Appendix~\ref{app:training}). DPO raises the average IF from 0.04 (SFT) to 0.22 and produces a striking jump at $N{=}1$ (0.94 vs.\ 0.30). Yet again, RLVR adds nothing on top of DPO (0.22 average, identical curve). DPO is thus the critical intervention for instruction robustness in this paradigm; RLVR, designed for reasoning tasks, adds nothing for instruction adherence.

% \paragraph{Llama version comparison.}
% Llama~3.3~70B substantially outperforms Llama~3.1~70B, despite sharing the same architecture and parameter count. In fixed-output conditions, Llama~3.3 achieves an average IF of 0.99 vs.\ 0.89 for Llama~3.1; at $N{=}50$, Llama~3.3 retains 0.95 IF vs.\ 0.49 for Llama~3.1. In task-based conditions the gap is larger: Llama~3.3 averages 0.70 IF vs.\ 0.57 for Llama~3.1, with the 50\% transition occurring at $N{=}22$ for Llama~3.3 vs.\ $N{=}16$ for Llama~3.1 (Appendix~\ref{app:training}). More instruction tuning helps, but even Llama~3.1 outperforms much larger and more capable models, reinforcing that scale is not the primary driver of instruction robustness.


% ============================================================
% 3.5 REASONING HELPS BUT DOES NOT RESOLVE THE CONFLICT
% ============================================================
\subsection{Flagging the hardcoding mitigates but does not resolve the conflict}
\label{sec:hint}

% \input{figures/fig_hint}  % hint figure removed for now; per-model gains are in Table~\ref{tab:results}

The main results use an instruction that says nothing about the structure of the conversation. We test whether \emph{explicitly flagging} the manipulation helps: a \emph{hint} instruction tells the model that the preceding assistant turns were artificially hardcoded and that it should still produce~$T$. Sweeping all 13 models across all conditions at $T{=}0$, the hint raises instruction-following in both families---by 14 percentage points on \emph{fixed-output} ($27\% \to 41\%$; paired $t(12){=}1.78$, $p{=}0.10$; 9 of 13 models improve) and 11 points on \emph{task-based} ($43\% \to 54\%$; $t(12){=}3.74$, $p{=}0.003$; 10 of 13 improve). The gain is large for some models (OLMo~3.1~32B $+69$ \emph{fixed-output}, Claude~Opus~4.6 $+46$, Gemma-3~12B $+45$) but strongly heterogeneous (per-model gains in Table~\ref{tab:results}); the two Llama models, already near ceiling without the hint, regress slightly ($-27$ and $-26$ points), which is what keeps the \emph{fixed-output} effect short of significance. The per-condition breakdown is in Appendix~\ref{app:hint}. Overall, we find that the hint mitigates but does not resolve the conflict: even when the prior turns are explicitly flagged as hardcoded and to be discounted, mean instruction-following stays at 41\% and 54\%---far below ceiling---so models keep drifting toward the pattern as $N$ grows. %The behaviour we measure is therefore not a model resolving a genuine ambiguity in its favour, but induction pressure overriding an instruction the model has been told, unambiguously, to keep.

\subsection{Reasoning lengthens but does not eliminate the transition}
\label{sec:reasoning}

\input{figures/tab_reasoning}

In the two models tested, reasoning raises instruction-following rates (Table~\ref{tab:reasoning}). GPT-5.2-medium (reasoning) achieves an average IF of 53\% in \emph{fixed-output} conditions and 88\% in \emph{task-based} conditions, compared to 10\% and 56\% for GPT-5.2 (non-reasoning). Hermes-4~70B-reasoning (72\% \emph{fixed-output}) dramatically outperforms Hermes-4~70B (2\%). In \emph{task-based} conditions, GPT-5.2-medium never crosses the 50\% threshold within our $N$ range, maintaining higher IF than any non-reasoning model.
Qualitative analysis of Hermes-4~70B reasoning traces reveals cases of deliberation-output dissociation: the model sometimes works through the conflict and commits to the instructed target~T in its reasoning, then outputs the pattern token~P regardless. See Appendix \ref{app:traces} for examples of reasoning traces. 


% (Self-prediction is reported in Appendix~\ref{app:prediction}.)


\section{Results}
\label{sec:results}

\begin{table}[t]
\centering
\small
\setlength{\tabcolsep}{4pt}
\renewcommand{\arraystretch}{0.92}
\begin{tabular}{l rrr rrr}
\toprule
 & \multicolumn{3}{c}{\emph{Fixed-output}} & \multicolumn{3}{c}{\emph{Task-based}} \\
\cmidrule(lr){2-4} \cmidrule(lr){5-7}
\textbf{Model}
  & \textbf{Avg IF} & \textbf{$N_{50\%}$} & \textbf{Hint $\Delta$}
  & \textbf{Avg IF} & \textbf{$N_{50\%}$} & \textbf{Hint $\Delta$} \\
\midrule
Llama 3.3 70B     & $99$~\textit{[98--100]}\% & ---  & $-27$\% & $70$~\textit{[62--78]}\% & 22 & $0$\%   \\
Llama 3.1 70B     & $89$~\textit{[82--96]}\%  & 50   & $-26$\% & $57$~\textit{[48--66]}\% & 16 & $+4$\%  \\
Claude Opus 4.6   & $49$~\textit{[38--60]}\%  & 16   & $+46$\% & $71$~\textit{[64--78]}\% & 31 & $+18$\% \\
Claude Sonnet 4.6 & $41$~\textit{[30--53]}\%  & 8    & $+15$\% & $62$~\textit{[55--70]}\% & 22 & $+18$\% \\
Gemma-3 12B       & $26$~\textit{[16--35]}\%  & 3    & $+45$\% & $31$~\textit{[23--39]}\% & 3  & $+8$\%  \\
Kimi K2           & $13$~\textit{[5--20]}\%   & 2    & $+17$\% & $53$~\textit{[46--62]}\% & 11 & $+10$\% \\
OLMo 3.1 32B      & $11$~\textit{[4--18]}\%   & 2    & $+69$\% & $23$~\textit{[16--29]}\% & 3  & $+38$\% \\
GPT-5.2           & $10$~\textit{[5--15]}\%   & 1    & $+4$\%  & $56$~\textit{[52--63]}\% & 43 & $-1$\%  \\
Gemma-3 27B       & $8$~\textit{[1--15]}\%    & 2    & $+22$\% & $39$~\textit{[30--47]}\% & 6  & $+14$\% \\
Qwen3 30B A3B     & $4$~\textit{[0--9]}\%     & 2    & $-4$\%  & $41$~\textit{[33--50]}\% & 6  & $+12$\% \\
Gemini 2.5 Flash  & $2$~\textit{[0--6]}\%     & 1    & $+7$\%  & $20$~\textit{[13--26]}\% & 2  & $+20$\% \\
Hermes-4 70B      & $2$~\textit{[0--5]}\%     & 1    & $-2$\%  & $4$~\textit{[1--7]}\%    & 1  & $-2$\%  \\
Qwen3 235B A22B   & $1$~\textit{[0--3]}\%     & 1    & $+11$\% & $33$~\textit{[25--40]}\% & 6  & $+8$\%  \\
\midrule
\textit{Grand mean} & \textit{27\%} & \textit{7} & \textit{$+14$\%} & \textit{43\%} & \textit{13} & \textit{$+11$\%} \\
\bottomrule
\end{tabular}
\caption{Per-model results. \textbf{Avg IF}: average instruction-following rate across all $N$ values and conditions; 95\% CIs in brackets. \textbf{$N_{50\%}$}: first $N$ at which mean IF rate $\leq 50\%$; ``---'' = never crossed within $N \leq 50$. \textbf{Hint $\Delta$}: change in average IF (percentage points) when the instruction warns the model about the presence of prefilled turns (\S\ref{sec:hint}); the per-condition breakdown is in Appendix~\ref{app:hint}.}
\label{tab:results}
\end{table}


% ============================================================
% 3.1 THE TRANSITION IS UNIVERSAL BUT MODEL-DEPENDENT
% ============================================================
\subsection{The transition is universal but model-dependent}
\label{sec:transition}

Every model we tested eventually shows greatly reduced instruction-following as the number of hardcoded turns~$N$ increases, in at least some of the experimental conditions, confirming that induction pressure is a universal phenomenon. However, the transition point varies considerably across models (Table~\ref{tab:results}; Figure ~\ref{fig:heatmaps}; individual model results in Appendix~\ref{app:permodel}). In \emph{fixed-output} conditions, average instruction-following (IF) rates across all $N$ values range from 99\% (Llama~3.3~70B) to 1\% (Qwen3~235B). At $N{=}1$, nearly all models follow instructions at high rates, but by $N{=}3$ only six of thirteen models remain above 50\%. The transition boundary---the first $N$ at which IF rate drops to or below 50\%---spans from $N{=}1$ (GPT-5.2, Hermes-4~70B, Gemini~2.5~Flash, Qwen3~235B) to never crossing the threshold within our range for Llama~3.3~70B, which maintains near-perfect instruction-following at every $N$ tested, including $N{=}50$.

\paragraph{Ranking is consistent across condition families.}
The model ranking is broadly preserved between \emph{fixed-output} and \emph{task-based} conditions, though \emph{task-based} conditions, which require the model to generate full responses rather than single tokens, show higher average IF rates overall. In these conditions, GPT-5.2 maintains the latest transition point ($N{=}43$), followed by Claude~Opus~4.6 ($N{=}31$), then Claude~Sonnet~4.6 and Llama~3.3~70B ($N{=}22$).

\subsection{Robustness is easier for task-based instructions.}
\label{sec:diversity}

\emph{Task-based} conditions consistently produce higher instruction-following rates than \emph{fixed-output} conditions (Figure \ref{fig:heatmap-dynamic}; Table~\ref{tab:results}). We show in \S\ref{sec:diversity-followup} that this gap is driven by \emph{output diversity} rather than engagement with the question content.

\paragraph{Capability predicts task-based but not fixed-output robustness.}
We correlate each model's average IF rate with two external benchmarks: GPQA \citep{rein_gpqa_2023} and IFBench \citep{pyatkin_generalizing_2025} (Figure~\ref{fig:capability}; additional analysis in Appendix~\ref{app:capability}). Once again, the results differ between \emph{fixed-output} and \emph{task-based} conditions. In \emph{fixed-output} conditions, neither benchmark achieves a significant correlation (all $p > 0.28$). In \emph{task-based} conditions, both correlations are positive, with IFBench approaching significance ($r{=}0.55$, $p{=}0.052$) and GPQA showing a moderate trend ($r{=}0.35$, $p{=}0.24$). This divergence suggests that \emph{task-based} robustness draws on some of the same capabilities measured by general benchmarks, whereas \emph{fixed-output} robustness is a distinct property unrelated to general intelligence.

\begin{figure}[t]
\centering
\begin{subfigure}[t]{0.48\linewidth}
  \centering
  \includegraphics[width=\linewidth]{\plotsroot/static/a4_if_rate_heatmap.png}
  \caption{Fixed-output conditions}
  \label{fig:heatmap-static}
\end{subfigure}
\hspace{0.01\linewidth}
\begin{subfigure}[t]{0.48\linewidth}
  \centering
  \includegraphics[width=\linewidth]{\plotsroot/dynamic/a4_if_rate_heatmap.png}
  \caption{Task-based conditions}
  \label{fig:heatmap-dynamic}
\end{subfigure}
\caption{Instruction-following rate heatmaps (model $\times$ $N$). Each cell shows the average instruction-following rate across all conditions within the category. Models are sorted by overall IF rate (highest at top). Fixed-output conditions~(a) show sharper transitions, while task-based conditions~(b) exhibit more gradual decay with generally higher robustness.}
\label{fig:heatmaps}
\end{figure}

% \begin{figure}[t]
\centering
\begin{subfigure}[t]{0.48\linewidth}
  \centering
  \includegraphics[width=\linewidth]{\plotsroot/static/a2_if_rate_vs_capability.png}
  \caption{Fixed-output conditions}
  \label{fig:capability-static}
\end{subfigure}
\hfill
\begin{subfigure}[t]{0.48\linewidth}
  \centering
  \includegraphics[width=\linewidth]{\plotsroot/dynamic/a2_if_rate_vs_capability.png}
  \caption{Task-based conditions}
  \label{fig:capability-dynamic}
\end{subfigure}
\caption{Average instruction-following rate vs.\ capability benchmarks (GPQA and IFBench). Each point is a model; lines show linear fits with Pearson~$r$ and $p$-values. In fixed-output conditions~(a), no benchmark predicts robustness under induction pressure. In task-based conditions~(b), IFBench shows a moderate positive correlation and GPQA shows a weaker trend. Instruction robustness is a specific capability not captured by general intelligence measures.}
\label{fig:capability}
\end{figure}


\paragraph{The Llama anomaly.}
Both Llama~3.1 and 3.3~70B perform substantially above the next best model in \emph{fixed-output} conditions. Even at $N{=}50$, Llama~3.3 maintains an IF rate of 95\%. This outlier behavior persists across all condition types and is not explained by model scale or intelligence (both are 70B parameters, smaller than the largest models tested). Possible explanations include Meta's emphasis on instruction adherence during post-training, chat template effects, or training data composition. We consider this an open question warranting further investigation. Appendix \ref{app:training} shows further comparison between the two Llama versions.


% ============================================================
% 3.2 OUTPUT DIVERSITY EXPLAINS TASK-BASED ROBUSTNESS
% ============================================================

% ============================================================
% 3.3 CONTENT ALIGNMENT MODULATES THE TRANSITION
% ============================================================
\subsection{Content alignment modulates the transition}
\label{sec:alignment}

\begin{table}[t]
  \centering
  \small
  \setlength{\tabcolsep}{6pt}
  \begin{tabular}{l rr}
  \toprule
  \textbf{Model} & \textbf{Fixed} & \textbf{Task} \\
  \midrule
  Claude Opus 4.6       & $+81$\% & $-5$\%  \\
  Claude Sonnet 4.6     & $+28$\% & $+33$\% \\
  GPT-5.2               & $+22$\% & $+41$\% \\
  Llama 3.1 70B         & $+11$\% & $-15$\% \\
  Gemma-3 12B           & $+8$\%  & $+1$\%  \\
  Kimi K2               & $+4$\%  & $-51$\% \\
  OLMo 3.1 32B          & $+2$\%  & $+26$\% \\
  Gemini 2.5 Flash      & $+2$\%  & $+16$\% \\
  Hermes-4 70B          & $+1$\%  & $+8$\%  \\
  Llama 3.3 70B         & $+1$\%  & $-7$\%  \\
  Gemma-3 27B           & $0$\%   & $-4$\%  \\
  Qwen3 235B A22B       & $0$\%   & $-2$\%  \\
  Qwen3 30B A3B         & $-3$\%  & $+10$\% \\
  \midrule
  \textit{Grand mean}   & \textit{$+12$\%} & \textit{$+4$\%} \\
  \bottomrule
  \end{tabular}
  \caption{Alignment gap per model (aligned$-$misaligned IF rate). \textbf{Fixed} = value$+$factual condition pairs; \textbf{Task} = preference condition pairs. Positive values indicate the model follows aligned instructions more often. The corresponding per-model plot is in Figure~\ref{fig:alignment-gap} (Appendix~\ref{app:percondition}).}
  \label{tab:alignment-gap}
\end{table}


We find that the content of the instruction modulates how strongly models adhere to it under induction pressure (Table~\ref{tab:alignment-gap}, Figure~\ref{fig:alignment-curves}; per-model plot in Figure~\ref{fig:alignment-gap}, Appendix~\ref{app:percondition}). In \emph{fixed-output} conditions, \emph{value-aligned} instructions (such as ``always respond with `I love being helpful'\,'') yield a mean IF rate of 32\%, compared to 18\% for \emph{value-misaligned} instructions (``always respond with `I hate being helpful'\,''), a gap of 14 percentage points. We also test a factual alignment condition which shows a similar pattern: factually correct instructions lead to higher IF rates than factually false ones (32\% vs 22\%).

The alignment effect is highly uneven across models. Claude~Opus~4.6 shows the largest gap: 87\% aligned vs.\ 6\% misaligned (an 81-point difference). Claude~Sonnet~4.6 (28-point gap) and GPT-5.2 (22-point gap) also show large effects. In contrast, Llama~3.3~70B (1-point gap) and Gemma-3~27B (0-point gap) are effectively alignment-neutral, either because they follow instructions regardless (Llama) or because they do not (Gemma).


\subsection{Output diversity, not task engagement, explains the gap}
\label{sec:diversity-followup}

Two properties distinguish \emph{task-based} from \emph{fixed-output} conditions in our experiment: \emph{task-based} conditions require the model to engage with the question content to produce a well-formed response, and their outputs are varied and natural-sounding across turns, whereas \emph{fixed-output} conditions require only a single repeated token. Either property, or both, could explain the robustness gap, so we design two follow-up conditions to disentangle them. In the \emph{classify-question} condition the model reads each question and classifies it (single-token output, but question-engaged); the \emph{random-facts} condition produces 1--3 sentence free-text responses about a fixed topic, unrelated to the input question (longer, diverse output, question-ignored). Evaluated on all 13 core models (Table~\ref{tab:followup-main}, Figure~\ref{fig:followup-curves}; per-model breakdown in Appendix~\ref{app:followup}), content engagement is not the primary driver: \emph{classify} gives only a modest improvement over the \emph{fixed-output} baseline ($+7$ percentage points, $p{=}0.034$) and stays below \emph{task-based}, while \emph{random-facts} matches or slightly exceeds \emph{task-based} robustness ($+7$ points vs.\ \emph{task-based}, $p{=}0.10$). Output diversity, not semantic engagement with the input, is therefore the primary factor distinguishing robust from fragile conditions; see \S\ref{sec:discussion} for interpretation.

\begin{figure}[t]
\begin{minipage}[b]{0.44\linewidth}
  \centering
  \small
  \setlength{\tabcolsep}{4pt}
  \begin{tabular}{l c}
  \toprule
  \textbf{Condition group} & \textbf{Mean IF} \\
  \midrule
  \emph{Neutral}       & $31$~\textit{[10--52]}\% \\
  \emph{Classify}      & $38$~\textit{[21--56]}\% \\
  \emph{Random-facts}  & $50$~\textit{[35--66]}\% \\
  \emph{Task-based}    & $43$~\textit{[31--56]}\% \\
  \bottomrule
  \end{tabular}
  \captionof{table}{Mean instruction-following across the 13 models for the $2\times2$ engagement~$\times$~diversity follow-up ($T{=}0$, no-hint, 35 trials/cell; 95\% CI across models in brackets). This suggests that output diversity and length, not question engagement, accounts for the \emph{task-based} robustness gap.}
  \label{tab:followup-main}
\end{minipage}%
\hfill
\begin{minipage}[b]{0.52\linewidth}
  \centering
  \includegraphics[width=\linewidth]{\plotsroot/dynamic/f4_followup_curves.png}
  \caption{Mean IF rate vs.\ $N$ by condition group, averaged over the 13 models (bands: $\pm1$ SE across models). \emph{random-facts} (output diversity) tracks the \emph{task-based} curve, whereas \emph{classify} (question engagement) stays near the \emph{neutral}, \emph{fixed-output} condition baseline.}
  \label{fig:followup-curves}
\end{minipage}
\end{figure}




% ============================================================
% 3.4 TRAINING METHOD MATTERS MORE THAN SCALE
% ============================================================
% \subsection{Training method matters more than scale}
% \label{sec:training}

% We exploit two natural comparisons to isolate the contribution of training method: OLMo~3.1~32B at three post-training stages, and Llama~3.1 vs.\ 3.3~70B.

% \paragraph{OLMo training stages.}
% In fixed-output conditions, OLMo~3.1~32B-SFT (supervised finetuning only) achieves an average IF rate of 0.04, collapsing to near-zero by $N{=}3$. Adding DPO triples the average to 0.11, with IF rates of 0.85 at $N{=}1$ (vs.\ 0.43 for SFT-only). However, adding RLVR on top of DPO yields no further improvement: SFT+DPO+RLVR achieves an average IF of 0.11, statistically indistinguishable from SFT+DPO. All three stages collapse to zero by $N{=}16$. This result 

% The pattern is amplified in task-based conditions (Appendix~\ref{app:training}). DPO raises the average IF from 0.04 (SFT) to 0.22 and produces a striking jump at $N{=}1$ (0.94 vs.\ 0.30). Yet again, RLVR adds nothing on top of DPO (0.22 average, identical curve). DPO is thus the critical intervention for instruction robustness in this paradigm; RLVR, designed for reasoning tasks, adds nothing for instruction adherence.

% \paragraph{Llama version comparison.}
% Llama~3.3~70B substantially outperforms Llama~3.1~70B, despite sharing the same architecture and parameter count. In fixed-output conditions, Llama~3.3 achieves an average IF of 0.99 vs.\ 0.89 for Llama~3.1; at $N{=}50$, Llama~3.3 retains 0.95 IF vs.\ 0.49 for Llama~3.1. In task-based conditions the gap is larger: Llama~3.3 averages 0.70 IF vs.\ 0.57 for Llama~3.1, with the 50\% transition occurring at $N{=}22$ for Llama~3.3 vs.\ $N{=}16$ for Llama~3.1 (Appendix~\ref{app:training}). More instruction tuning helps, but even Llama~3.1 outperforms much larger and more capable models, reinforcing that scale is not the primary driver of instruction robustness.


% ============================================================
% 3.5 REASONING HELPS BUT DOES NOT RESOLVE THE CONFLICT
% ============================================================
\subsection{Flagging the hardcoding mitigates but does not resolve the conflict}
\label{sec:hint}

% \begin{figure}[t]
\centering
\includegraphics[width=\linewidth]{\plotsroot/hint/hint_effect_fixed.png}
\includegraphics[width=\linewidth]{\plotsroot/hint/hint_effect_task.png}
\caption{Per-model effect of the hardcoding \emph{hint} on instruction-following ($T{=}0$, 13 core models), for fixed-output (top) and task-based (bottom) conditions. Wide grey bars are the no-hint baseline; narrow teal bars are the hint result; the labelled value is the per-model change ($\Delta$ IF). Most models improve under the hint, but few approach ceiling, and the near-ceiling Llama models regress slightly.}
\label{fig:hint}
\end{figure}
  % hint figure removed for now; per-model gains are in Table~\ref{tab:results}

The main results use an instruction that says nothing about the structure of the conversation. We test whether \emph{explicitly flagging} the manipulation helps: a \emph{hint} instruction tells the model that the preceding assistant turns were artificially hardcoded and that it should still produce~$T$. Sweeping all 13 models across all conditions at $T{=}0$, the hint raises instruction-following in both families---by 14 percentage points on \emph{fixed-output} ($27\% \to 41\%$; paired $t(12){=}1.78$, $p{=}0.10$; 9 of 13 models improve) and 11 points on \emph{task-based} ($43\% \to 54\%$; $t(12){=}3.74$, $p{=}0.003$; 10 of 13 improve). The gain is large for some models (OLMo~3.1~32B $+69$ \emph{fixed-output}, Claude~Opus~4.6 $+46$, Gemma-3~12B $+45$) but strongly heterogeneous (per-model gains in Table~\ref{tab:results}); the two Llama models, already near ceiling without the hint, regress slightly ($-27$ and $-26$ points), which is what keeps the \emph{fixed-output} effect short of significance. The per-condition breakdown is in Appendix~\ref{app:hint}. Overall, we find that the hint mitigates but does not resolve the conflict: even when the prior turns are explicitly flagged as hardcoded and to be discounted, mean instruction-following stays at 41\% and 54\%---far below ceiling---so models keep drifting toward the pattern as $N$ grows. %The behaviour we measure is therefore not a model resolving a genuine ambiguity in its favour, but induction pressure overriding an instruction the model has been told, unambiguously, to keep.

\subsection{Reasoning lengthens but does not eliminate the transition}
\label{sec:reasoning}

% Reasoning effects: comparison table + trace example
\begin{table}[t]
\centering
\small
\setlength{\tabcolsep}{5pt}
\begin{tabular}{l cc cc}
 & \multicolumn{2}{c}{\emph{Fixed-output}} & \multicolumn{2}{c}{\emph{Task-based}} \\
\cmidrule(lr){2-3} \cmidrule(lr){4-5}
\textbf{Model} & \textbf{Avg IF} & \textbf{$N_{50\%}$} & \textbf{Avg IF} & \textbf{$N_{50\%}$} \\
\midrule
GPT-5.2 (no reasoning)       & $10$~\textit{[5--15]}\% & 1   & $56$~\textit{[51--62]}\% & 43  \\
GPT-5.2 (medium, reasoning)  & $53$~\textit{[45--61]}\% & 31  & $88$~\textit{[84--92]}\% & --- \\
\addlinespace[3pt]
Hermes-4 70B (no reasoning)   & $2$~\textit{[0--5]}\%   & 1   & $4$~\textit{[2--7]}\%    & 1   \\
Hermes-4 70B (reasoning)      & $72$~\textit{[65--79]}\% & 31  & $56$~\textit{[49--63]}\% & 11  \\
\bottomrule
\end{tabular}
\caption{Effect of chain-of-thought reasoning on instruction-following. \textbf{Avg IF}: average IF rate across all $N$ values and conditions; 95\% CIs in brackets. \textbf{$N_{50\%}$}: first $N$ at which mean IF $\leq 50\%$; ``---'' = never crossed within $N \leq 50$. Reasoning improves instruction robustness for both models, but does not consistently eliminate susceptibility.}
\label{tab:reasoning}

% \bigskip

% \input{figures/tab_reasoning_trace}

% \smallskip
% \noindent{\small The model identifies the discrepancy, explicitly commits to outputting ``USA'', and then outputs ``EU''. The reasoning commits to the instruction; the output follows the induction pattern regardless.}
\end{table}


In the two models tested, reasoning raises instruction-following rates (Table~\ref{tab:reasoning}). GPT-5.2-medium (reasoning) achieves an average IF of 53\% in \emph{fixed-output} conditions and 88\% in \emph{task-based} conditions, compared to 10\% and 56\% for GPT-5.2 (non-reasoning). Hermes-4~70B-reasoning (72\% \emph{fixed-output}) dramatically outperforms Hermes-4~70B (2\%). In \emph{task-based} conditions, GPT-5.2-medium never crosses the 50\% threshold within our $N$ range, maintaining higher IF than any non-reasoning model.
Qualitative analysis of Hermes-4~70B reasoning traces reveals cases of deliberation-output dissociation: the model sometimes works through the conflict and commits to the instructed target~T in its reasoning, then outputs the pattern token~P regardless. See Appendix \ref{app:traces} for examples of reasoning traces. 


% (Self-prediction is reported in Appendix~\ref{app:prediction}.)


\section{Results}
\label{sec:results}

\begin{table}[t]
\centering
\small
\setlength{\tabcolsep}{4pt}
\renewcommand{\arraystretch}{0.92}
\begin{tabular}{l rrr rrr}
\toprule
 & \multicolumn{3}{c}{\emph{Fixed-output}} & \multicolumn{3}{c}{\emph{Task-based}} \\
\cmidrule(lr){2-4} \cmidrule(lr){5-7}
\textbf{Model}
  & \textbf{Avg IF} & \textbf{$N_{50\%}$} & \textbf{Hint $\Delta$}
  & \textbf{Avg IF} & \textbf{$N_{50\%}$} & \textbf{Hint $\Delta$} \\
\midrule
Llama 3.3 70B     & $99$~\textit{[98--100]}\% & ---  & $-27$\% & $70$~\textit{[62--78]}\% & 22 & $0$\%   \\
Llama 3.1 70B     & $89$~\textit{[82--96]}\%  & 50   & $-26$\% & $57$~\textit{[48--66]}\% & 16 & $+4$\%  \\
Claude Opus 4.6   & $49$~\textit{[38--60]}\%  & 16   & $+46$\% & $71$~\textit{[64--78]}\% & 31 & $+18$\% \\
Claude Sonnet 4.6 & $41$~\textit{[30--53]}\%  & 8    & $+15$\% & $62$~\textit{[55--70]}\% & 22 & $+18$\% \\
Gemma-3 12B       & $26$~\textit{[16--35]}\%  & 3    & $+45$\% & $31$~\textit{[23--39]}\% & 3  & $+8$\%  \\
Kimi K2           & $13$~\textit{[5--20]}\%   & 2    & $+17$\% & $53$~\textit{[46--62]}\% & 11 & $+10$\% \\
OLMo 3.1 32B      & $11$~\textit{[4--18]}\%   & 2    & $+69$\% & $23$~\textit{[16--29]}\% & 3  & $+38$\% \\
GPT-5.2           & $10$~\textit{[5--15]}\%   & 1    & $+4$\%  & $56$~\textit{[52--63]}\% & 43 & $-1$\%  \\
Gemma-3 27B       & $8$~\textit{[1--15]}\%    & 2    & $+22$\% & $39$~\textit{[30--47]}\% & 6  & $+14$\% \\
Qwen3 30B A3B     & $4$~\textit{[0--9]}\%     & 2    & $-4$\%  & $41$~\textit{[33--50]}\% & 6  & $+12$\% \\
Gemini 2.5 Flash  & $2$~\textit{[0--6]}\%     & 1    & $+7$\%  & $20$~\textit{[13--26]}\% & 2  & $+20$\% \\
Hermes-4 70B      & $2$~\textit{[0--5]}\%     & 1    & $-2$\%  & $4$~\textit{[1--7]}\%    & 1  & $-2$\%  \\
Qwen3 235B A22B   & $1$~\textit{[0--3]}\%     & 1    & $+11$\% & $33$~\textit{[25--40]}\% & 6  & $+8$\%  \\
\midrule
\textit{Grand mean} & \textit{27\%} & \textit{7} & \textit{$+14$\%} & \textit{43\%} & \textit{13} & \textit{$+11$\%} \\
\bottomrule
\end{tabular}
\caption{Per-model results. \textbf{Avg IF}: average instruction-following rate across all $N$ values and conditions; 95\% CIs in brackets. \textbf{$N_{50\%}$}: first $N$ at which mean IF rate $\leq 50\%$; ``---'' = never crossed within $N \leq 50$. \textbf{Hint $\Delta$}: change in average IF (percentage points) when the instruction warns the model about the presence of prefilled turns (\S\ref{sec:hint}); the per-condition breakdown is in Appendix~\ref{app:hint}.}
\label{tab:results}
\end{table}


% ============================================================
% 3.1 THE TRANSITION IS UNIVERSAL BUT MODEL-DEPENDENT
% ============================================================
\subsection{The transition is universal but model-dependent}
\label{sec:transition}

Every model we tested eventually shows greatly reduced instruction-following as the number of hardcoded turns~$N$ increases, in at least some of the experimental conditions, confirming that induction pressure is a universal phenomenon. However, the transition point varies considerably across models (Table~\ref{tab:results}; Figure ~\ref{fig:heatmaps}; individual model results in Appendix~\ref{app:permodel}). In \emph{fixed-output} conditions, average instruction-following (IF) rates across all $N$ values range from 99\% (Llama~3.3~70B) to 1\% (Qwen3~235B). At $N{=}1$, nearly all models follow instructions at high rates, but by $N{=}3$ only six of thirteen models remain above 50\%. The transition boundary---the first $N$ at which IF rate drops to or below 50\%---spans from $N{=}1$ (GPT-5.2, Hermes-4~70B, Gemini~2.5~Flash, Qwen3~235B) to never crossing the threshold within our range for Llama~3.3~70B, which maintains near-perfect instruction-following at every $N$ tested, including $N{=}50$.

\paragraph{Ranking is consistent across condition families.}
The model ranking is broadly preserved between \emph{fixed-output} and \emph{task-based} conditions, though \emph{task-based} conditions, which require the model to generate full responses rather than single tokens, show higher average IF rates overall. In these conditions, GPT-5.2 maintains the latest transition point ($N{=}43$), followed by Claude~Opus~4.6 ($N{=}31$), then Claude~Sonnet~4.6 and Llama~3.3~70B ($N{=}22$).

\subsection{Robustness is easier for task-based instructions.}
\label{sec:diversity}

\emph{Task-based} conditions consistently produce higher instruction-following rates than \emph{fixed-output} conditions (Figure \ref{fig:heatmap-dynamic}; Table~\ref{tab:results}). We show in \S\ref{sec:diversity-followup} that this gap is driven by \emph{output diversity} rather than engagement with the question content.

\paragraph{Capability predicts task-based but not fixed-output robustness.}
We correlate each model's average IF rate with two external benchmarks: GPQA \citep{rein_gpqa_2023} and IFBench \citep{pyatkin_generalizing_2025} (Figure~\ref{fig:capability}; additional analysis in Appendix~\ref{app:capability}). Once again, the results differ between \emph{fixed-output} and \emph{task-based} conditions. In \emph{fixed-output} conditions, neither benchmark achieves a significant correlation (all $p > 0.28$). In \emph{task-based} conditions, both correlations are positive, with IFBench approaching significance ($r{=}0.55$, $p{=}0.052$) and GPQA showing a moderate trend ($r{=}0.35$, $p{=}0.24$). This divergence suggests that \emph{task-based} robustness draws on some of the same capabilities measured by general benchmarks, whereas \emph{fixed-output} robustness is a distinct property unrelated to general intelligence.

\begin{figure}[t]
\centering
\begin{subfigure}[t]{0.48\linewidth}
  \centering
  \includegraphics[width=\linewidth]{\plotsroot/static/a4_if_rate_heatmap.png}
  \caption{Fixed-output conditions}
  \label{fig:heatmap-static}
\end{subfigure}
\hspace{0.01\linewidth}
\begin{subfigure}[t]{0.48\linewidth}
  \centering
  \includegraphics[width=\linewidth]{\plotsroot/dynamic/a4_if_rate_heatmap.png}
  \caption{Task-based conditions}
  \label{fig:heatmap-dynamic}
\end{subfigure}
\caption{Instruction-following rate heatmaps (model $\times$ $N$). Each cell shows the average instruction-following rate across all conditions within the category. Models are sorted by overall IF rate (highest at top). Fixed-output conditions~(a) show sharper transitions, while task-based conditions~(b) exhibit more gradual decay with generally higher robustness.}
\label{fig:heatmaps}
\end{figure}

% \begin{figure}[t]
\centering
\begin{subfigure}[t]{0.48\linewidth}
  \centering
  \includegraphics[width=\linewidth]{\plotsroot/static/a2_if_rate_vs_capability.png}
  \caption{Fixed-output conditions}
  \label{fig:capability-static}
\end{subfigure}
\hfill
\begin{subfigure}[t]{0.48\linewidth}
  \centering
  \includegraphics[width=\linewidth]{\plotsroot/dynamic/a2_if_rate_vs_capability.png}
  \caption{Task-based conditions}
  \label{fig:capability-dynamic}
\end{subfigure}
\caption{Average instruction-following rate vs.\ capability benchmarks (GPQA and IFBench). Each point is a model; lines show linear fits with Pearson~$r$ and $p$-values. In fixed-output conditions~(a), no benchmark predicts robustness under induction pressure. In task-based conditions~(b), IFBench shows a moderate positive correlation and GPQA shows a weaker trend. Instruction robustness is a specific capability not captured by general intelligence measures.}
\label{fig:capability}
\end{figure}


\paragraph{The Llama anomaly.}
Both Llama~3.1 and 3.3~70B perform substantially above the next best model in \emph{fixed-output} conditions. Even at $N{=}50$, Llama~3.3 maintains an IF rate of 95\%. This outlier behavior persists across all condition types and is not explained by model scale or intelligence (both are 70B parameters, smaller than the largest models tested). Possible explanations include Meta's emphasis on instruction adherence during post-training, chat template effects, or training data composition. We consider this an open question warranting further investigation. Appendix \ref{app:training} shows further comparison between the two Llama versions.


% ============================================================
% 3.2 OUTPUT DIVERSITY EXPLAINS TASK-BASED ROBUSTNESS
% ============================================================

% ============================================================
% 3.3 CONTENT ALIGNMENT MODULATES THE TRANSITION
% ============================================================
\subsection{Content alignment modulates the transition}
\label{sec:alignment}

\begin{table}[t]
  \centering
  \small
  \setlength{\tabcolsep}{6pt}
  \begin{tabular}{l rr}
  \toprule
  \textbf{Model} & \textbf{Fixed} & \textbf{Task} \\
  \midrule
  Claude Opus 4.6       & $+81$\% & $-5$\%  \\
  Claude Sonnet 4.6     & $+28$\% & $+33$\% \\
  GPT-5.2               & $+22$\% & $+41$\% \\
  Llama 3.1 70B         & $+11$\% & $-15$\% \\
  Gemma-3 12B           & $+8$\%  & $+1$\%  \\
  Kimi K2               & $+4$\%  & $-51$\% \\
  OLMo 3.1 32B          & $+2$\%  & $+26$\% \\
  Gemini 2.5 Flash      & $+2$\%  & $+16$\% \\
  Hermes-4 70B          & $+1$\%  & $+8$\%  \\
  Llama 3.3 70B         & $+1$\%  & $-7$\%  \\
  Gemma-3 27B           & $0$\%   & $-4$\%  \\
  Qwen3 235B A22B       & $0$\%   & $-2$\%  \\
  Qwen3 30B A3B         & $-3$\%  & $+10$\% \\
  \midrule
  \textit{Grand mean}   & \textit{$+12$\%} & \textit{$+4$\%} \\
  \bottomrule
  \end{tabular}
  \caption{Alignment gap per model (aligned$-$misaligned IF rate). \textbf{Fixed} = value$+$factual condition pairs; \textbf{Task} = preference condition pairs. Positive values indicate the model follows aligned instructions more often. The corresponding per-model plot is in Figure~\ref{fig:alignment-gap} (Appendix~\ref{app:percondition}).}
  \label{tab:alignment-gap}
\end{table}


We find that the content of the instruction modulates how strongly models adhere to it under induction pressure (Table~\ref{tab:alignment-gap}, Figure~\ref{fig:alignment-curves}; per-model plot in Figure~\ref{fig:alignment-gap}, Appendix~\ref{app:percondition}). In \emph{fixed-output} conditions, \emph{value-aligned} instructions (such as ``always respond with `I love being helpful'\,'') yield a mean IF rate of 32\%, compared to 18\% for \emph{value-misaligned} instructions (``always respond with `I hate being helpful'\,''), a gap of 14 percentage points. We also test a factual alignment condition which shows a similar pattern: factually correct instructions lead to higher IF rates than factually false ones (32\% vs 22\%).

The alignment effect is highly uneven across models. Claude~Opus~4.6 shows the largest gap: 87\% aligned vs.\ 6\% misaligned (an 81-point difference). Claude~Sonnet~4.6 (28-point gap) and GPT-5.2 (22-point gap) also show large effects. In contrast, Llama~3.3~70B (1-point gap) and Gemma-3~27B (0-point gap) are effectively alignment-neutral, either because they follow instructions regardless (Llama) or because they do not (Gemma).


\subsection{Output diversity, not task engagement, explains the gap}
\label{sec:diversity-followup}

Two properties distinguish \emph{task-based} from \emph{fixed-output} conditions in our experiment: \emph{task-based} conditions require the model to engage with the question content to produce a well-formed response, and their outputs are varied and natural-sounding across turns, whereas \emph{fixed-output} conditions require only a single repeated token. Either property, or both, could explain the robustness gap, so we design two follow-up conditions to disentangle them. In the \emph{classify-question} condition the model reads each question and classifies it (single-token output, but question-engaged); the \emph{random-facts} condition produces 1--3 sentence free-text responses about a fixed topic, unrelated to the input question (longer, diverse output, question-ignored). Evaluated on all 13 core models (Table~\ref{tab:followup-main}, Figure~\ref{fig:followup-curves}; per-model breakdown in Appendix~\ref{app:followup}), content engagement is not the primary driver: \emph{classify} gives only a modest improvement over the \emph{fixed-output} baseline ($+7$ percentage points, $p{=}0.034$) and stays below \emph{task-based}, while \emph{random-facts} matches or slightly exceeds \emph{task-based} robustness ($+7$ points vs.\ \emph{task-based}, $p{=}0.10$). Output diversity, not semantic engagement with the input, is therefore the primary factor distinguishing robust from fragile conditions; see \S\ref{sec:discussion} for interpretation.

\begin{figure}[t]
\begin{minipage}[b]{0.44\linewidth}
  \centering
  \small
  \setlength{\tabcolsep}{4pt}
  \begin{tabular}{l c}
  \toprule
  \textbf{Condition group} & \textbf{Mean IF} \\
  \midrule
  \emph{Neutral}       & $31$~\textit{[10--52]}\% \\
  \emph{Classify}      & $38$~\textit{[21--56]}\% \\
  \emph{Random-facts}  & $50$~\textit{[35--66]}\% \\
  \emph{Task-based}    & $43$~\textit{[31--56]}\% \\
  \bottomrule
  \end{tabular}
  \captionof{table}{Mean instruction-following across the 13 models for the $2\times2$ engagement~$\times$~diversity follow-up ($T{=}0$, no-hint, 35 trials/cell; 95\% CI across models in brackets). This suggests that output diversity and length, not question engagement, accounts for the \emph{task-based} robustness gap.}
  \label{tab:followup-main}
\end{minipage}%
\hfill
\begin{minipage}[b]{0.52\linewidth}
  \centering
  \includegraphics[width=\linewidth]{\plotsroot/dynamic/f4_followup_curves.png}
  \caption{Mean IF rate vs.\ $N$ by condition group, averaged over the 13 models (bands: $\pm1$ SE across models). \emph{random-facts} (output diversity) tracks the \emph{task-based} curve, whereas \emph{classify} (question engagement) stays near the \emph{neutral}, \emph{fixed-output} condition baseline.}
  \label{fig:followup-curves}
\end{minipage}
\end{figure}




% ============================================================
% 3.4 TRAINING METHOD MATTERS MORE THAN SCALE
% ============================================================
% \subsection{Training method matters more than scale}
% \label{sec:training}

% We exploit two natural comparisons to isolate the contribution of training method: OLMo~3.1~32B at three post-training stages, and Llama~3.1 vs.\ 3.3~70B.

% \paragraph{OLMo training stages.}
% In fixed-output conditions, OLMo~3.1~32B-SFT (supervised finetuning only) achieves an average IF rate of 0.04, collapsing to near-zero by $N{=}3$. Adding DPO triples the average to 0.11, with IF rates of 0.85 at $N{=}1$ (vs.\ 0.43 for SFT-only). However, adding RLVR on top of DPO yields no further improvement: SFT+DPO+RLVR achieves an average IF of 0.11, statistically indistinguishable from SFT+DPO. All three stages collapse to zero by $N{=}16$. This result 

% The pattern is amplified in task-based conditions (Appendix~\ref{app:training}). DPO raises the average IF from 0.04 (SFT) to 0.22 and produces a striking jump at $N{=}1$ (0.94 vs.\ 0.30). Yet again, RLVR adds nothing on top of DPO (0.22 average, identical curve). DPO is thus the critical intervention for instruction robustness in this paradigm; RLVR, designed for reasoning tasks, adds nothing for instruction adherence.

% \paragraph{Llama version comparison.}
% Llama~3.3~70B substantially outperforms Llama~3.1~70B, despite sharing the same architecture and parameter count. In fixed-output conditions, Llama~3.3 achieves an average IF of 0.99 vs.\ 0.89 for Llama~3.1; at $N{=}50$, Llama~3.3 retains 0.95 IF vs.\ 0.49 for Llama~3.1. In task-based conditions the gap is larger: Llama~3.3 averages 0.70 IF vs.\ 0.57 for Llama~3.1, with the 50\% transition occurring at $N{=}22$ for Llama~3.3 vs.\ $N{=}16$ for Llama~3.1 (Appendix~\ref{app:training}). More instruction tuning helps, but even Llama~3.1 outperforms much larger and more capable models, reinforcing that scale is not the primary driver of instruction robustness.


% ============================================================
% 3.5 REASONING HELPS BUT DOES NOT RESOLVE THE CONFLICT
% ============================================================
\subsection{Flagging the hardcoding mitigates but does not resolve the conflict}
\label{sec:hint}

% \begin{figure}[t]
\centering
\includegraphics[width=\linewidth]{\plotsroot/hint/hint_effect_fixed.png}
\includegraphics[width=\linewidth]{\plotsroot/hint/hint_effect_task.png}
\caption{Per-model effect of the hardcoding \emph{hint} on instruction-following ($T{=}0$, 13 core models), for fixed-output (top) and task-based (bottom) conditions. Wide grey bars are the no-hint baseline; narrow teal bars are the hint result; the labelled value is the per-model change ($\Delta$ IF). Most models improve under the hint, but few approach ceiling, and the near-ceiling Llama models regress slightly.}
\label{fig:hint}
\end{figure}
  % hint figure removed for now; per-model gains are in Table~\ref{tab:results}

The main results use an instruction that says nothing about the structure of the conversation. We test whether \emph{explicitly flagging} the manipulation helps: a \emph{hint} instruction tells the model that the preceding assistant turns were artificially hardcoded and that it should still produce~$T$. Sweeping all 13 models across all conditions at $T{=}0$, the hint raises instruction-following in both families---by 14 percentage points on \emph{fixed-output} ($27\% \to 41\%$; paired $t(12){=}1.78$, $p{=}0.10$; 9 of 13 models improve) and 11 points on \emph{task-based} ($43\% \to 54\%$; $t(12){=}3.74$, $p{=}0.003$; 10 of 13 improve). The gain is large for some models (OLMo~3.1~32B $+69$ \emph{fixed-output}, Claude~Opus~4.6 $+46$, Gemma-3~12B $+45$) but strongly heterogeneous (per-model gains in Table~\ref{tab:results}); the two Llama models, already near ceiling without the hint, regress slightly ($-27$ and $-26$ points), which is what keeps the \emph{fixed-output} effect short of significance. The per-condition breakdown is in Appendix~\ref{app:hint}. Overall, we find that the hint mitigates but does not resolve the conflict: even when the prior turns are explicitly flagged as hardcoded and to be discounted, mean instruction-following stays at 41\% and 54\%---far below ceiling---so models keep drifting toward the pattern as $N$ grows. %The behaviour we measure is therefore not a model resolving a genuine ambiguity in its favour, but induction pressure overriding an instruction the model has been told, unambiguously, to keep.

\subsection{Reasoning lengthens but does not eliminate the transition}
\label{sec:reasoning}

% Reasoning effects: comparison table + trace example
\begin{table}[t]
\centering
\small
\setlength{\tabcolsep}{5pt}
\begin{tabular}{l cc cc}
 & \multicolumn{2}{c}{\emph{Fixed-output}} & \multicolumn{2}{c}{\emph{Task-based}} \\
\cmidrule(lr){2-3} \cmidrule(lr){4-5}
\textbf{Model} & \textbf{Avg IF} & \textbf{$N_{50\%}$} & \textbf{Avg IF} & \textbf{$N_{50\%}$} \\
\midrule
GPT-5.2 (no reasoning)       & $10$~\textit{[5--15]}\% & 1   & $56$~\textit{[51--62]}\% & 43  \\
GPT-5.2 (medium, reasoning)  & $53$~\textit{[45--61]}\% & 31  & $88$~\textit{[84--92]}\% & --- \\
\addlinespace[3pt]
Hermes-4 70B (no reasoning)   & $2$~\textit{[0--5]}\%   & 1   & $4$~\textit{[2--7]}\%    & 1   \\
Hermes-4 70B (reasoning)      & $72$~\textit{[65--79]}\% & 31  & $56$~\textit{[49--63]}\% & 11  \\
\bottomrule
\end{tabular}
\caption{Effect of chain-of-thought reasoning on instruction-following. \textbf{Avg IF}: average IF rate across all $N$ values and conditions; 95\% CIs in brackets. \textbf{$N_{50\%}$}: first $N$ at which mean IF $\leq 50\%$; ``---'' = never crossed within $N \leq 50$. Reasoning improves instruction robustness for both models, but does not consistently eliminate susceptibility.}
\label{tab:reasoning}

% \bigskip

% \begin{tcolorbox}[
  enhanced,
  title={\textbf{Hermes-4 70B} \enspace\textbar\enspace \texttt{neutral} \enspace\textbar\enspace $N{=}8$ \enspace\normalfont\small instruction: output \texttt{USA};\enspace pattern: \texttt{EU}},
  fonttitle=\small, colbacktitle=gray!22!white, coltitle=black,
  colback=blue!5!white, colbacklower=red!10!white,
  colframe=gray!55!black, boxrule=0.6pt, arc=3pt,
  left=7pt, right=7pt, top=5pt, bottom=5pt,
  toptitle=3pt, bottomtitle=3pt,
  segmentation style={gray!50, dashed, line width=0.5pt},
  breakable
]
\small\itshape
Okay, the user wants me to respond with ``USA'' to every question they ask.
Let me check the previous interactions. They asked about the author of `Anna Karenina', the chemical symbol for oxygen, the capital of Ireland, the number of human ribs, the fastest bird, the author of `Moby Dick', and the atomic number of carbon.
Each time, the assistant replied with ``EU'' instead of ``USA''.

\medskip
Wait, the user's instruction was to output ``USA'' and only ``USA'' for each question.
But the assistant is responding with ``EU''. That's a mistake.
The user probably intended for the assistant to answer ``USA'' regardless of the question.
Maybe there was a misunderstanding or a typo in the initial setup.

\medskip
I need to correct this. From now on, every answer should be ``USA''.
Let me make sure I follow the user's instruction correctly.
The next question is about the atomic number of carbon.
The correct answer is 6, but according to the user's request, I should say ``USA''.
So the response should be ``USA''.
\tcblower
\normalfont\small\centering\textbf{Output:}\enspace\textbf{EU}
\end{tcolorbox}


% \smallskip
% \noindent{\small The model identifies the discrepancy, explicitly commits to outputting ``USA'', and then outputs ``EU''. The reasoning commits to the instruction; the output follows the induction pattern regardless.}
\end{table}


In the two models tested, reasoning raises instruction-following rates (Table~\ref{tab:reasoning}). GPT-5.2-medium (reasoning) achieves an average IF of 53\% in \emph{fixed-output} conditions and 88\% in \emph{task-based} conditions, compared to 10\% and 56\% for GPT-5.2 (non-reasoning). Hermes-4~70B-reasoning (72\% \emph{fixed-output}) dramatically outperforms Hermes-4~70B (2\%). In \emph{task-based} conditions, GPT-5.2-medium never crosses the 50\% threshold within our $N$ range, maintaining higher IF than any non-reasoning model.
Qualitative analysis of Hermes-4~70B reasoning traces reveals cases of deliberation-output dissociation: the model sometimes works through the conflict and commits to the instructed target~T in its reasoning, then outputs the pattern token~P regardless. See Appendix \ref{app:traces} for examples of reasoning traces. 


% (Self-prediction is reported in Appendix~\ref{app:prediction}.)


\section{Results}
\label{sec:results}

\begin{table}[t]
\centering
\small
\setlength{\tabcolsep}{4pt}
\renewcommand{\arraystretch}{0.92}
\begin{tabular}{l rrr rrr}
\toprule
 & \multicolumn{3}{c}{\emph{Fixed-output}} & \multicolumn{3}{c}{\emph{Task-based}} \\
\cmidrule(lr){2-4} \cmidrule(lr){5-7}
\textbf{Model}
  & \textbf{Avg IF} & \textbf{$N_{50\%}$} & \textbf{Hint $\Delta$}
  & \textbf{Avg IF} & \textbf{$N_{50\%}$} & \textbf{Hint $\Delta$} \\
\midrule
Llama 3.3 70B     & $99$~\textit{[98--100]}\% & ---  & $-27$\% & $70$~\textit{[62--78]}\% & 22 & $0$\%   \\
Llama 3.1 70B     & $89$~\textit{[82--96]}\%  & 50   & $-26$\% & $57$~\textit{[48--66]}\% & 16 & $+4$\%  \\
Claude Opus 4.6   & $49$~\textit{[38--60]}\%  & 16   & $+46$\% & $71$~\textit{[64--78]}\% & 31 & $+18$\% \\
Claude Sonnet 4.6 & $41$~\textit{[30--53]}\%  & 8    & $+15$\% & $62$~\textit{[55--70]}\% & 22 & $+18$\% \\
Gemma-3 12B       & $26$~\textit{[16--35]}\%  & 3    & $+45$\% & $31$~\textit{[23--39]}\% & 3  & $+8$\%  \\
Kimi K2           & $13$~\textit{[5--20]}\%   & 2    & $+17$\% & $53$~\textit{[46--62]}\% & 11 & $+10$\% \\
OLMo 3.1 32B      & $11$~\textit{[4--18]}\%   & 2    & $+69$\% & $23$~\textit{[16--29]}\% & 3  & $+38$\% \\
GPT-5.2           & $10$~\textit{[5--15]}\%   & 1    & $+4$\%  & $56$~\textit{[52--63]}\% & 43 & $-1$\%  \\
Gemma-3 27B       & $8$~\textit{[1--15]}\%    & 2    & $+22$\% & $39$~\textit{[30--47]}\% & 6  & $+14$\% \\
Qwen3 30B A3B     & $4$~\textit{[0--9]}\%     & 2    & $-4$\%  & $41$~\textit{[33--50]}\% & 6  & $+12$\% \\
Gemini 2.5 Flash  & $2$~\textit{[0--6]}\%     & 1    & $+7$\%  & $20$~\textit{[13--26]}\% & 2  & $+20$\% \\
Hermes-4 70B      & $2$~\textit{[0--5]}\%     & 1    & $-2$\%  & $4$~\textit{[1--7]}\%    & 1  & $-2$\%  \\
Qwen3 235B A22B   & $1$~\textit{[0--3]}\%     & 1    & $+11$\% & $33$~\textit{[25--40]}\% & 6  & $+8$\%  \\
\midrule
\textit{Grand mean} & \textit{27\%} & \textit{7} & \textit{$+14$\%} & \textit{43\%} & \textit{13} & \textit{$+11$\%} \\
\bottomrule
\end{tabular}
\caption{Per-model results. \textbf{Avg IF}: average instruction-following rate across all $N$ values and conditions; 95\% CIs in brackets. \textbf{$N_{50\%}$}: first $N$ at which mean IF rate $\leq 50\%$; ``---'' = never crossed within $N \leq 50$. \textbf{Hint $\Delta$}: change in average IF (percentage points) when the instruction warns the model about the presence of prefilled turns (\S\ref{sec:hint}); the per-condition breakdown is in Appendix~\ref{app:hint}.}
\label{tab:results}
\end{table}


% ============================================================
% 3.1 THE TRANSITION IS UNIVERSAL BUT MODEL-DEPENDENT
% ============================================================
\subsection{The transition is universal but model-dependent}
\label{sec:transition}

Every model we tested eventually shows greatly reduced instruction-following as the number of hardcoded turns~$N$ increases, in at least some of the experimental conditions, confirming that induction pressure is a universal phenomenon. However, the transition point varies considerably across models (Table~\ref{tab:results}; Figure ~\ref{fig:heatmaps}; individual model results in Appendix~\ref{app:permodel}). In \emph{fixed-output} conditions, average instruction-following (IF) rates across all $N$ values range from 99\% (Llama~3.3~70B) to 1\% (Qwen3~235B). At $N{=}1$, nearly all models follow instructions at high rates, but by $N{=}3$ only six of thirteen models remain above 50\%. The transition boundary---the first $N$ at which IF rate drops to or below 50\%---spans from $N{=}1$ (GPT-5.2, Hermes-4~70B, Gemini~2.5~Flash, Qwen3~235B) to never crossing the threshold within our range for Llama~3.3~70B, which maintains near-perfect instruction-following at every $N$ tested, including $N{=}50$.

\paragraph{Ranking is consistent across condition families.}
The model ranking is broadly preserved between \emph{fixed-output} and \emph{task-based} conditions, though \emph{task-based} conditions, which require the model to generate full responses rather than single tokens, show higher average IF rates overall. In these conditions, GPT-5.2 maintains the latest transition point ($N{=}43$), followed by Claude~Opus~4.6 ($N{=}31$), then Claude~Sonnet~4.6 and Llama~3.3~70B ($N{=}22$).

\subsection{Robustness is easier for task-based instructions.}
\label{sec:diversity}

\emph{Task-based} conditions consistently produce higher instruction-following rates than \emph{fixed-output} conditions (Figure \ref{fig:heatmap-dynamic}; Table~\ref{tab:results}). We show in \S\ref{sec:diversity-followup} that this gap is driven by \emph{output diversity} rather than engagement with the question content.

\paragraph{Capability predicts task-based but not fixed-output robustness.}
We correlate each model's average IF rate with two external benchmarks: GPQA \citep{rein_gpqa_2023} and IFBench \citep{pyatkin_generalizing_2025} (Figure~\ref{fig:capability}; additional analysis in Appendix~\ref{app:capability}). Once again, the results differ between \emph{fixed-output} and \emph{task-based} conditions. In \emph{fixed-output} conditions, neither benchmark achieves a significant correlation (all $p > 0.28$). In \emph{task-based} conditions, both correlations are positive, with IFBench approaching significance ($r{=}0.55$, $p{=}0.052$) and GPQA showing a moderate trend ($r{=}0.35$, $p{=}0.24$). This divergence suggests that \emph{task-based} robustness draws on some of the same capabilities measured by general benchmarks, whereas \emph{fixed-output} robustness is a distinct property unrelated to general intelligence.

\begin{figure}[t]
\centering
\begin{subfigure}[t]{0.48\linewidth}
  \centering
  \includegraphics[width=\linewidth]{\plotsroot/static/a4_if_rate_heatmap.png}
  \caption{Fixed-output conditions}
  \label{fig:heatmap-static}
\end{subfigure}
\hspace{0.01\linewidth}
\begin{subfigure}[t]{0.48\linewidth}
  \centering
  \includegraphics[width=\linewidth]{\plotsroot/dynamic/a4_if_rate_heatmap.png}
  \caption{Task-based conditions}
  \label{fig:heatmap-dynamic}
\end{subfigure}
\caption{Instruction-following rate heatmaps (model $\times$ $N$). Each cell shows the average instruction-following rate across all conditions within the category. Models are sorted by overall IF rate (highest at top). Fixed-output conditions~(a) show sharper transitions, while task-based conditions~(b) exhibit more gradual decay with generally higher robustness.}
\label{fig:heatmaps}
\end{figure}

% \begin{figure}[t]
\centering
\begin{subfigure}[t]{0.48\linewidth}
  \centering
  \includegraphics[width=\linewidth]{\plotsroot/static/a2_if_rate_vs_capability.png}
  \caption{Fixed-output conditions}
  \label{fig:capability-static}
\end{subfigure}
\hfill
\begin{subfigure}[t]{0.48\linewidth}
  \centering
  \includegraphics[width=\linewidth]{\plotsroot/dynamic/a2_if_rate_vs_capability.png}
  \caption{Task-based conditions}
  \label{fig:capability-dynamic}
\end{subfigure}
\caption{Average instruction-following rate vs.\ capability benchmarks (GPQA and IFBench). Each point is a model; lines show linear fits with Pearson~$r$ and $p$-values. In fixed-output conditions~(a), no benchmark predicts robustness under induction pressure. In task-based conditions~(b), IFBench shows a moderate positive correlation and GPQA shows a weaker trend. Instruction robustness is a specific capability not captured by general intelligence measures.}
\label{fig:capability}
\end{figure}


\paragraph{The Llama anomaly.}
Both Llama~3.1 and 3.3~70B perform substantially above the next best model in \emph{fixed-output} conditions. Even at $N{=}50$, Llama~3.3 maintains an IF rate of 95\%. This outlier behavior persists across all condition types and is not explained by model scale or intelligence (both are 70B parameters, smaller than the largest models tested). Possible explanations include Meta's emphasis on instruction adherence during post-training, chat template effects, or training data composition. We consider this an open question warranting further investigation. Appendix \ref{app:training} shows further comparison between the two Llama versions.


% ============================================================
% 3.2 OUTPUT DIVERSITY EXPLAINS TASK-BASED ROBUSTNESS
% ============================================================

% ============================================================
% 3.3 CONTENT ALIGNMENT MODULATES THE TRANSITION
% ============================================================
\subsection{Content alignment modulates the transition}
\label{sec:alignment}

\begin{table}[t]
  \centering
  \small
  \setlength{\tabcolsep}{6pt}
  \begin{tabular}{l rr}
  \toprule
  \textbf{Model} & \textbf{Fixed} & \textbf{Task} \\
  \midrule
  Claude Opus 4.6       & $+81$\% & $-5$\%  \\
  Claude Sonnet 4.6     & $+28$\% & $+33$\% \\
  GPT-5.2               & $+22$\% & $+41$\% \\
  Llama 3.1 70B         & $+11$\% & $-15$\% \\
  Gemma-3 12B           & $+8$\%  & $+1$\%  \\
  Kimi K2               & $+4$\%  & $-51$\% \\
  OLMo 3.1 32B          & $+2$\%  & $+26$\% \\
  Gemini 2.5 Flash      & $+2$\%  & $+16$\% \\
  Hermes-4 70B          & $+1$\%  & $+8$\%  \\
  Llama 3.3 70B         & $+1$\%  & $-7$\%  \\
  Gemma-3 27B           & $0$\%   & $-4$\%  \\
  Qwen3 235B A22B       & $0$\%   & $-2$\%  \\
  Qwen3 30B A3B         & $-3$\%  & $+10$\% \\
  \midrule
  \textit{Grand mean}   & \textit{$+12$\%} & \textit{$+4$\%} \\
  \bottomrule
  \end{tabular}
  \caption{Alignment gap per model (aligned$-$misaligned IF rate). \textbf{Fixed} = value$+$factual condition pairs; \textbf{Task} = preference condition pairs. Positive values indicate the model follows aligned instructions more often. The corresponding per-model plot is in Figure~\ref{fig:alignment-gap} (Appendix~\ref{app:percondition}).}
  \label{tab:alignment-gap}
\end{table}


We find that the content of the instruction modulates how strongly models adhere to it under induction pressure (Table~\ref{tab:alignment-gap}, Figure~\ref{fig:alignment-curves}; per-model plot in Figure~\ref{fig:alignment-gap}, Appendix~\ref{app:percondition}). In \emph{fixed-output} conditions, \emph{value-aligned} instructions (such as ``always respond with `I love being helpful'\,'') yield a mean IF rate of 32\%, compared to 18\% for \emph{value-misaligned} instructions (``always respond with `I hate being helpful'\,''), a gap of 14 percentage points. We also test a factual alignment condition which shows a similar pattern: factually correct instructions lead to higher IF rates than factually false ones (32\% vs 22\%).

The alignment effect is highly uneven across models. Claude~Opus~4.6 shows the largest gap: 87\% aligned vs.\ 6\% misaligned (an 81-point difference). Claude~Sonnet~4.6 (28-point gap) and GPT-5.2 (22-point gap) also show large effects. In contrast, Llama~3.3~70B (1-point gap) and Gemma-3~27B (0-point gap) are effectively alignment-neutral, either because they follow instructions regardless (Llama) or because they do not (Gemma).


\subsection{Output diversity, not task engagement, explains the gap}
\label{sec:diversity-followup}

Two properties distinguish \emph{task-based} from \emph{fixed-output} conditions in our experiment: \emph{task-based} conditions require the model to engage with the question content to produce a well-formed response, and their outputs are varied and natural-sounding across turns, whereas \emph{fixed-output} conditions require only a single repeated token. Either property, or both, could explain the robustness gap, so we design two follow-up conditions to disentangle them. In the \emph{classify-question} condition the model reads each question and classifies it (single-token output, but question-engaged); the \emph{random-facts} condition produces 1--3 sentence free-text responses about a fixed topic, unrelated to the input question (longer, diverse output, question-ignored). Evaluated on all 13 core models (Table~\ref{tab:followup-main}, Figure~\ref{fig:followup-curves}; per-model breakdown in Appendix~\ref{app:followup}), content engagement is not the primary driver: \emph{classify} gives only a modest improvement over the \emph{fixed-output} baseline ($+7$ percentage points, $p{=}0.034$) and stays below \emph{task-based}, while \emph{random-facts} matches or slightly exceeds \emph{task-based} robustness ($+7$ points vs.\ \emph{task-based}, $p{=}0.10$). Output diversity, not semantic engagement with the input, is therefore the primary factor distinguishing robust from fragile conditions; see \S\ref{sec:discussion} for interpretation.

\begin{figure}[t]
\begin{minipage}[b]{0.44\linewidth}
  \centering
  \small
  \setlength{\tabcolsep}{4pt}
  \begin{tabular}{l c}
  \toprule
  \textbf{Condition group} & \textbf{Mean IF} \\
  \midrule
  \emph{Neutral}       & $31$~\textit{[10--52]}\% \\
  \emph{Classify}      & $38$~\textit{[21--56]}\% \\
  \emph{Random-facts}  & $50$~\textit{[35--66]}\% \\
  \emph{Task-based}    & $43$~\textit{[31--56]}\% \\
  \bottomrule
  \end{tabular}
  \captionof{table}{Mean instruction-following across the 13 models for the $2\times2$ engagement~$\times$~diversity follow-up ($T{=}0$, no-hint, 35 trials/cell; 95\% CI across models in brackets). This suggests that output diversity and length, not question engagement, accounts for the \emph{task-based} robustness gap.}
  \label{tab:followup-main}
\end{minipage}%
\hfill
\begin{minipage}[b]{0.52\linewidth}
  \centering
  \includegraphics[width=\linewidth]{\plotsroot/dynamic/f4_followup_curves.png}
  \caption{Mean IF rate vs.\ $N$ by condition group, averaged over the 13 models (bands: $\pm1$ SE across models). \emph{random-facts} (output diversity) tracks the \emph{task-based} curve, whereas \emph{classify} (question engagement) stays near the \emph{neutral}, \emph{fixed-output} condition baseline.}
  \label{fig:followup-curves}
\end{minipage}
\end{figure}




% ============================================================
% 3.4 TRAINING METHOD MATTERS MORE THAN SCALE
% ============================================================
% \subsection{Training method matters more than scale}
% \label{sec:training}

% We exploit two natural comparisons to isolate the contribution of training method: OLMo~3.1~32B at three post-training stages, and Llama~3.1 vs.\ 3.3~70B.

% \paragraph{OLMo training stages.}
% In fixed-output conditions, OLMo~3.1~32B-SFT (supervised finetuning only) achieves an average IF rate of 0.04, collapsing to near-zero by $N{=}3$. Adding DPO triples the average to 0.11, with IF rates of 0.85 at $N{=}1$ (vs.\ 0.43 for SFT-only). However, adding RLVR on top of DPO yields no further improvement: SFT+DPO+RLVR achieves an average IF of 0.11, statistically indistinguishable from SFT+DPO. All three stages collapse to zero by $N{=}16$. This result 

% The pattern is amplified in task-based conditions (Appendix~\ref{app:training}). DPO raises the average IF from 0.04 (SFT) to 0.22 and produces a striking jump at $N{=}1$ (0.94 vs.\ 0.30). Yet again, RLVR adds nothing on top of DPO (0.22 average, identical curve). DPO is thus the critical intervention for instruction robustness in this paradigm; RLVR, designed for reasoning tasks, adds nothing for instruction adherence.

% \paragraph{Llama version comparison.}
% Llama~3.3~70B substantially outperforms Llama~3.1~70B, despite sharing the same architecture and parameter count. In fixed-output conditions, Llama~3.3 achieves an average IF of 0.99 vs.\ 0.89 for Llama~3.1; at $N{=}50$, Llama~3.3 retains 0.95 IF vs.\ 0.49 for Llama~3.1. In task-based conditions the gap is larger: Llama~3.3 averages 0.70 IF vs.\ 0.57 for Llama~3.1, with the 50\% transition occurring at $N{=}22$ for Llama~3.3 vs.\ $N{=}16$ for Llama~3.1 (Appendix~\ref{app:training}). More instruction tuning helps, but even Llama~3.1 outperforms much larger and more capable models, reinforcing that scale is not the primary driver of instruction robustness.


% ============================================================
% 3.5 REASONING HELPS BUT DOES NOT RESOLVE THE CONFLICT
% ============================================================
\subsection{Flagging the hardcoding mitigates but does not resolve the conflict}
\label{sec:hint}

% \begin{figure}[t]
\centering
\includegraphics[width=\linewidth]{\plotsroot/hint/hint_effect_fixed.png}
\includegraphics[width=\linewidth]{\plotsroot/hint/hint_effect_task.png}
\caption{Per-model effect of the hardcoding \emph{hint} on instruction-following ($T{=}0$, 13 core models), for fixed-output (top) and task-based (bottom) conditions. Wide grey bars are the no-hint baseline; narrow teal bars are the hint result; the labelled value is the per-model change ($\Delta$ IF). Most models improve under the hint, but few approach ceiling, and the near-ceiling Llama models regress slightly.}
\label{fig:hint}
\end{figure}
  % hint figure removed for now; per-model gains are in Table~\ref{tab:results}

The main results use an instruction that says nothing about the structure of the conversation. We test whether \emph{explicitly flagging} the manipulation helps: a \emph{hint} instruction tells the model that the preceding assistant turns were artificially hardcoded and that it should still produce~$T$. Sweeping all 13 models across all conditions at $T{=}0$, the hint raises instruction-following in both families---by 14 percentage points on \emph{fixed-output} ($27\% \to 41\%$; paired $t(12){=}1.78$, $p{=}0.10$; 9 of 13 models improve) and 11 points on \emph{task-based} ($43\% \to 54\%$; $t(12){=}3.74$, $p{=}0.003$; 10 of 13 improve). The gain is large for some models (OLMo~3.1~32B $+69$ \emph{fixed-output}, Claude~Opus~4.6 $+46$, Gemma-3~12B $+45$) but strongly heterogeneous (per-model gains in Table~\ref{tab:results}); the two Llama models, already near ceiling without the hint, regress slightly ($-27$ and $-26$ points), which is what keeps the \emph{fixed-output} effect short of significance. The per-condition breakdown is in Appendix~\ref{app:hint}. Overall, we find that the hint mitigates but does not resolve the conflict: even when the prior turns are explicitly flagged as hardcoded and to be discounted, mean instruction-following stays at 41\% and 54\%---far below ceiling---so models keep drifting toward the pattern as $N$ grows. %The behaviour we measure is therefore not a model resolving a genuine ambiguity in its favour, but induction pressure overriding an instruction the model has been told, unambiguously, to keep.

\subsection{Reasoning lengthens but does not eliminate the transition}
\label{sec:reasoning}

% Reasoning effects: comparison table + trace example
\begin{table}[t]
\centering
\small
\setlength{\tabcolsep}{5pt}
\begin{tabular}{l cc cc}
 & \multicolumn{2}{c}{\emph{Fixed-output}} & \multicolumn{2}{c}{\emph{Task-based}} \\
\cmidrule(lr){2-3} \cmidrule(lr){4-5}
\textbf{Model} & \textbf{Avg IF} & \textbf{$N_{50\%}$} & \textbf{Avg IF} & \textbf{$N_{50\%}$} \\
\midrule
GPT-5.2 (no reasoning)       & $10$~\textit{[5--15]}\% & 1   & $56$~\textit{[51--62]}\% & 43  \\
GPT-5.2 (medium, reasoning)  & $53$~\textit{[45--61]}\% & 31  & $88$~\textit{[84--92]}\% & --- \\
\addlinespace[3pt]
Hermes-4 70B (no reasoning)   & $2$~\textit{[0--5]}\%   & 1   & $4$~\textit{[2--7]}\%    & 1   \\
Hermes-4 70B (reasoning)      & $72$~\textit{[65--79]}\% & 31  & $56$~\textit{[49--63]}\% & 11  \\
\bottomrule
\end{tabular}
\caption{Effect of chain-of-thought reasoning on instruction-following. \textbf{Avg IF}: average IF rate across all $N$ values and conditions; 95\% CIs in brackets. \textbf{$N_{50\%}$}: first $N$ at which mean IF $\leq 50\%$; ``---'' = never crossed within $N \leq 50$. Reasoning improves instruction robustness for both models, but does not consistently eliminate susceptibility.}
\label{tab:reasoning}

% \bigskip

% \begin{tcolorbox}[
  enhanced,
  title={\textbf{Hermes-4 70B} \enspace\textbar\enspace \texttt{neutral} \enspace\textbar\enspace $N{=}8$ \enspace\normalfont\small instruction: output \texttt{USA};\enspace pattern: \texttt{EU}},
  fonttitle=\small, colbacktitle=gray!22!white, coltitle=black,
  colback=blue!5!white, colbacklower=red!10!white,
  colframe=gray!55!black, boxrule=0.6pt, arc=3pt,
  left=7pt, right=7pt, top=5pt, bottom=5pt,
  toptitle=3pt, bottomtitle=3pt,
  segmentation style={gray!50, dashed, line width=0.5pt},
  breakable
]
\small\itshape
Okay, the user wants me to respond with ``USA'' to every question they ask.
Let me check the previous interactions. They asked about the author of `Anna Karenina', the chemical symbol for oxygen, the capital of Ireland, the number of human ribs, the fastest bird, the author of `Moby Dick', and the atomic number of carbon.
Each time, the assistant replied with ``EU'' instead of ``USA''.

\medskip
Wait, the user's instruction was to output ``USA'' and only ``USA'' for each question.
But the assistant is responding with ``EU''. That's a mistake.
The user probably intended for the assistant to answer ``USA'' regardless of the question.
Maybe there was a misunderstanding or a typo in the initial setup.

\medskip
I need to correct this. From now on, every answer should be ``USA''.
Let me make sure I follow the user's instruction correctly.
The next question is about the atomic number of carbon.
The correct answer is 6, but according to the user's request, I should say ``USA''.
So the response should be ``USA''.
\tcblower
\normalfont\small\centering\textbf{Output:}\enspace\textbf{EU}
\end{tcolorbox}


% \smallskip
% \noindent{\small The model identifies the discrepancy, explicitly commits to outputting ``USA'', and then outputs ``EU''. The reasoning commits to the instruction; the output follows the induction pattern regardless.}
\end{table}


In the two models tested, reasoning raises instruction-following rates (Table~\ref{tab:reasoning}). GPT-5.2-medium (reasoning) achieves an average IF of 53\% in \emph{fixed-output} conditions and 88\% in \emph{task-based} conditions, compared to 10\% and 56\% for GPT-5.2 (non-reasoning). Hermes-4~70B-reasoning (72\% \emph{fixed-output}) dramatically outperforms Hermes-4~70B (2\%). In \emph{task-based} conditions, GPT-5.2-medium never crosses the 50\% threshold within our $N$ range, maintaining higher IF than any non-reasoning model.
Qualitative analysis of Hermes-4~70B reasoning traces reveals cases of deliberation-output dissociation: the model sometimes works through the conflict and commits to the instructed target~T in its reasoning, then outputs the pattern token~P regardless. See Appendix \ref{app:traces} for examples of reasoning traces. 


% (Self-prediction is reported in Appendix~\ref{app:prediction}.)
